% ============================================================
% PRÉ-PROJETO DE TCC - EVASÃO PROEJA IFG GOIÂNIA
% Compilar com: xelatex + biber (ou latexmk)
% ============================================================
\documentclass[12pt,a4paper]{article}

% ---- PACOTES ESSENCIAIS ----
\usepackage{enumitem} % Permite usar \item[Label] em itemize
\usepackage[top=3cm, bottom=2cm, left=3cm, right=2cm]{geometry}
\usepackage[brazil]{babel}
\usepackage[utf8]{inputenc}
\usepackage[T1]{fontenc}
\usepackage{setspace}
\usepackage{indentfirst} % Recuo no primeiro parágrafo
\usepackage[backend=biber,style=abnt,extrayear]{biblatex}
\usepackage{hyperref}
\usepackage{graphicx}
\usepackage{booktabs}
\usepackage{titlesec}
\usepackage{tocloft}
\usepackage{fancyhdr}
\usepackage{abntex2cite}
% Definição do comando \fonte
\newcommand{\fonte}[1]{\textbf{Fonte:} #1}

% ---- CONFIGURAÇÕES ABNT ----
\onehalfspacing % Espaçamento 1,5
\setlength{\parindent}{1.25cm} % Recuo de parágrafo

% Formatação de Títulos (ABNT: Destaque progressivo)
\titleformat{\section}{\bfseries\uppercase}{\thesection}{1em}{}
\titleformat{\subsection}{\bfseries}{\thesubsection}{1em}{}
\titleformat{\subsubsection}{\bfseries\itshape}{\thesubsubsection}{1em}{}

% Configuração de Sumário
\renewcommand{\contentsname}{SUMÁRIO}
\renewcommand{\cftsecfont}{\bfseries}

% ---- REFERÊNCIAS ----
\addbibresource{refs.bib}

\begin{document}

% ===== CAPA =====
\begin{titlepage}
  \begin{center}
    \uppercase{Instituto Federal de Educação, Ciência e Tecnologia de Goiás}\\
    \uppercase{Campus Goiânia}\\
    Departamento das Áreas Acadêmicas IV\\
    Coordenação de Informática\\
    Bacharelado em Sistemas de Informação\\[2.5cm]

    \textbf{PAULO BRANQUINHO DE SOUZA}\\[3.5cm]

    \textbf{\uppercase{Análise de evasão no PROEJA do IFG Campus Goiânia: identificação de padrões, propostas de intervenção e desenvolvimento de um dashboard para monitoramento}}\\[3.5cm]

    Pré-projeto de Trabalho de Conclusão de Curso apresentado ao Instituto Federal de Goiás como requisito básico para a conclusão do curso de Bacharelado em Sistemas de Informação.\\[1cm]

    Orientador: Prof. <título>. <nome>\\[3cm]

    Goiânia/GO\\
    2026
  \end{center}
\end{titlepage}

% ===== FICHA CATALOGRÁFICA (Exemplo de layout) =====
\newpage
\thispagestyle{empty}
\vspace*{\fill}
\begin{center}
\fbox{\begin{minipage}[c]{12cm}
\small
\vspace{0.5cm}
\hspace{0.5cm} Souza, Paulo Branquinho de. \\
\hspace{0.5cm} Análise de evasão no PROEJA do IFG Campus Goiânia / Paulo Branquinho de Souza. -- 2026. \\
\hspace{0.5cm} 20 p. \\
\vspace{0.3cm}
\hspace{0.5cm} Orientador: Prof. <nome>. \\
\hspace{0.5cm} Pré-projeto (Bacharelado) -- Instituto Federal de Goiás, 2026. \\
\vspace{0.3cm}
\hspace{0.5cm} 1. Evasão escolar. 2. EJA. 3. Dashboard. I. Título. \\
\vspace{0.5cm}
\end{minipage}}
\end{center}

% ===== SUMÁRIO =====
\newpage
\tableofcontents
\newpage

% ---- HIPERLINKS ----
\hypersetup{
  colorlinks=true,
  linkcolor=black,
  urlcolor=blue,
  citecolor=black
}

% ---- ARQUIVO DE REFERÊNCIAS ----
\begin{filecontents}[overwrite]{refs.bib}
@book{freire2023,
  author    = {Paulo Freire},
  title     = {Pedagogia do Oprimido},
  edition   = {87},
  publisher = {Paz e Terra},
  address   = {Rio de Janeiro},
  year      = {2023}
}
@article{dipierro2005,
  author    = {Maria Clara Di Pierro},
  title     = {Notas sobre a redefinição da identidade e das políticas públicas da educação de jovens e adultos no Brasil},
  journal   = {Educação \& Sociedade},
  volume    = {26},
  number    = {92},
  pages     = {1115--1139},
  year      = {2005}
}
@book{arroyo2017,
  author    = {Miguel Arroyo},
  title     = {Passageiros da noite: do trabalho para a EJA, itinerários pelo direito a uma vida justa},
  publisher = {Vozes},
  address   = {Petrópolis},
  year      = {2017}
}
@article{oliveira2021,
  author    = {Paula Lucas Oliveira and Nilva Celestina Do Carmo},
  title     = {A temática evasão escolar no contexto do PROEJA: uma revisão integrativa},
  journal   = {Revista Ponto de Vista},
  volume    = {10},
  number    = {1},
  pages     = {01--21},
  year      = {2021}
}
@article{nery2020,
  author    = {Maria Clara Ramos Nery and Juliana Oliveira Pereira and Patricia Guerreiro Do Amaral},
  title     = {Educação de Jovens e Adultos — EJA — seus impasses e suas contradições},
  journal   = {Revista Eletrônica Científica da UERGS},
  volume    = {6},
  number    = {1},
  pages     = {29--41},
  year      = {2020}
}
@article{castro2016,
  author    = {Mad'Ana Desirée Ribeiro De Castro and Jacqueline Maria Barbosa Vitorette},
  title     = {O Programa de Integração da Educação Profissional com a Educação Básica na Modalidade de Educação de Jovens e Adultos (PROEJA) no IFG - Câmpus Goiânia: um percurso contraditório na construção do direito à educação},
  journal   = {HOLOS},
  volume    = {2},
  pages     = {301--318},
  year      = {2016}
}
@mastersthesis{pereira2011,
  author    = {Josué Vidal Pereira},
  title     = {O PROEJA no Instituto Federal de Goiás -- Campus Goiânia: um estudo sobre os fatores de acesso e permanência na escola},
  school    = {Universidade de Brasília},
  address   = {Brasília},
  year      = {2011}
}
@mastersthesis{castro2018,
  author    = {Mad'Ana Desirée Ribeiro De Castro},
  title     = {O PROEJA no IFG - Câmpus Goiânia: contradições, limites e perspectivas},
  school    = {Pontifícia Universidade Católica de Goiás},
  address   = {Goiânia},
  year      = {2018}
}
@article{santos2025,
  author    = {Áurea Fernandes Dos Santos and Salete De Almeida Lima Brigato and Everson Carlos Marin Lopes and Luana De Araújo Almeida},
  title     = {Desafios da Educação de Jovens e Adultos no Brasil: entre metodologias inadequadas e políticas instáveis},
  journal   = {Missioneira},
  volume    = {27},
  number    = {1},
  pages     = {103--114},
  year      = {2025}
}
@article{dutra2026,
  author    = {Maria de Fátima Calógeras Dutra and Andréa Carla de Lima Melo and Rozineide Iraci Pereira da Silva},
  title     = {Desafios e Perspectivas da Educação de Jovens e Adultos (EJA): um olhar sobre a evasão e a prática pedagógica no Brasil Contemporâneo},
  journal   = {REASE},
  volume    = {12},
  number    = {5},
  year      = {2026}
}
@book{tinto1975,
  author    = {Vincent Tinto},
  title     = {Dropout from Higher Education: A Theoretical Synthesis of Recent Research},
  journal   = {Review of Educational Research},
  volume    = {45},
  number    = {1},
  pages     = {89--125},
  year      = {1975}
}
@book{frignotto2005,
  author    = {Gaudêncio Frigotto and Maria Ciavatta and Marise Ramos},
  title     = {Ensino médio integrado: concepção e contradições},
  publisher = {Cortez},
  address   = {São Paulo},
  year      = {2005}
}
@article{assis2017,
  author    = {Lucas Rocha Soares De Assis},
  title     = {Perfil de evasão no ensino superior brasileiro: uma abordagem de mineração de dados},
  school    = {Universidade de Brasília},
  address   = {Brasília},
  year      = {2017}
}
@article{silva2023,
  author    = {Tânia das Graças de Castro e Silva and Rosângela de Bessa Barbosa da Silva and Gislene Lisboa de Oliveira},
  title     = {Educação de Jovens e Adultos em Goiás e Goiânia à luz dos dados do INEP: balanço analítico das metas 8, 9 e 10},
  journal   = {Revista Sapiência},
  volume    = {12},
  number    = {2},
  pages     = {88--104},
  year      = {2023}
}
@book{gil2008,
  author    = {Antônio Carlos Gil},
  title     = {Métodos e técnicas de pesquisa social},
  edition   = {6},
  publisher = {Atlas},
  address   = {São Paulo},
  year      = {2008}
}
@book{creswell2010,
  author    = {John W. Creswell},
  title     = {Projeto de pesquisa: métodos qualitativo, quantitativo e misto},
  edition   = {3},
  publisher = {Artmed},
  address   = {Porto Alegre},
  year      = {2010}
}
@book{tolentino2025,
  author    = {Dorgival Tolentino Filho},
  title     = {Educação de jovens e adultos no Brasil: Avanços, desafios e perspectivas},
  journal   = {International Integralize Scientific},
  volume    = {5},
  number    = {46},
  year      = {2025}
}
@techreport{todos2025,
  author    = {{Todos pela Educação}},
  title     = {Anuário Brasileiro da Educação Básica 2025},
  address   = {São Paulo},
  year      = {2025}
}
@misc{inep2024,
  author    = {{INEP}},
  title     = {Censo Escolar da Educação Básica 2024},
  address   = {Brasília},
  year      = {2024}
}
@article{baker2014,
  author    = {Ryan S. Baker and Paul S. Inventado},
  title     = {Educational Data Mining and Learning Analytics},
  booktitle = {Learning Analytics},
  publisher = {Springer},
  pages     = {61--75},
  year      = {2014}
}
@book{few2006,
  author    = {Stephen Few},
  title     = {Information Dashboard Design: The Effective Visual Communication of Data},
  publisher = {O'Reilly Media},
  year      = {2006}
}
@book{mckinney2018,
  author    = {Wes McKinney},
  title     = {Python for Data Analysis},
  edition   = {2},
  publisher = {O'Reilly Media},
  year      = {2018}
}
@book{prodanov2013,
  author    = {Cleber C. Prodanov and Ernani C. de Freitas},
  title     = {Metodologia do trabalho científico: métodos e técnicas da pesquisa e do trabalho acadêmico},
  edition   = {2},
  publisher = {Feevale},
  address   = {Novo Hamburgo},
  year      = {2013}
}
@book{hastie2009,
  author    = {Trevor Hastie and Robert Tibshirani and Jerome Friedman},
  title     = {The Elements of Statistical Learning},
  edition   = {2},
  publisher = {Springer},
  year      = {2009}
}
@book{tufte2001,
  author    = {Edward R. Tufte},
  title     = {The Visual Display of Quantitative Information},
  edition   = {2},
  publisher = {Graphics Press},
  year      = {2001}
}
@misc{mec2005,
  author    = {{MEC}},
  title     = {Decreto nº 5.478, de 24 de junho de 2005},
  year      = {2005}
}
@misc{mec2006,
  author    = {{MEC}},
  title     = {Decreto nº 5.840, de 13 de julho de 2006},
  year      = {2006}
}
@inproceedings{machado2011,
  author    = {Maria Margarida Machado},
  title     = {A educação de jovens e adultos no Brasil após a Lei nº 10.639/2003 e a Lei nº 11.645/2008: um olhar sobre as pesquisas},
  booktitle = {Reunião Nacional da ANPED},
  year      = {2011},
  address   = {Natal}
}
@book{ferrao2003,
  author    = {Maria Eugénia Ferrão and Maria João Valente Rosa},
  title     = {Os determinantes do abandono escolar no ensino superior},
  publisher = {Fundação Calouste Gulbenkian},
  address   = {Lisboa},
  year      = {2003}
}
@article{silva2019,
  author    = {Edilene Pereira da Silva and Marcilia Chaves Barreto},
  title     = {Evasão escolar na Educação de Jovens e Adultos: fatores que influenciam o abandono},
  journal   = {Revista Educação e Emancipação},
  volume    = {12},
  number    = {3},
  pages     = {127--148},
  year      = {2019}
}
@article{behrens2019,
  author    = {Marilda Aparecida Behrens and Patrícia Osório},
  title     = {Docência na EJA: desafios e possibilidades},
  journal   = {Revista Diálogo Educacional},
  volume    = {19},
  number    = {61},
  pages     = {823--842},
  year      = {2019}
}
\end{filecontents}

% ---- INÍCIO DO DOCUMENTO ----
\begin{document}

% ===== CAPA =====
\begin{titlepage}
  \begin{center}
    \textbf{INSTITUTO FEDERAL DE EDUCAÇÃO, CIÊNCIA E TECNOLOGIA DE GOIÁS}\\[0.2cm]
    \textbf{CAMPUS GOIÂNIA}\\[0.2cm]
    DEPARTAMENTO DAS ÁREAS ACADÊMICAS IV\\[0.2cm]
    COORDENAÇÃO DE INFORMÁTICA\\[0.2cm]
    \textbf{BACHARELADO EM SISTEMAS DE INFORMAÇÃO}\\[2.5cm]

    \textbf{PAULO BRANQUINHO DE SOUZA}\\[3.5cm]

    \textbf{ANÁLISE DE EVASÃO NO PROEJA DO IFG CAMPUS GOIÂNIA:}\\[0.15cm]
    \textbf{IDENTIFICAÇÃO DE PADRÕES, PROPOSTAS DE INTERVENÇÃO}\\[0.15cm]
    \textbf{E DESENVOLVIMENTO DE UM DASHBOARD PARA MONITORAMENTO}\\[3.5cm]

    PRÉ-PROJETO DE TRABALHO DE CONCLUSÃO DE CURSO\\[0.3cm]
    \onehalfspacing
    Pré-projeto de Trabalho de Conclusão de Curso apresentado no Instituto Federal de Goiás como requisito básico para a conclusão do curso de Bacharelado em Sistemas de Informação.\\[1cm]

    Orientador: Prof. \textit{<título>}. \textit{<nome>}\\[3cm]

    Goiânia/GO\\
    2026
  \end{center}
\end{titlepage}

% ===== FICHA CATALOGRÁFICA =====
\newpage
\thispagestyle{empty}
\begin{center}
  \textbf{Paulo Branquinho de Souza}
\end{center}

\noindent Análise de evasão no PROEJA do IFG Campus Goiânia: identificação de padrões, propostas de intervenção e desenvolvimento de um dashboard para monitoramento / Paulo Branquinho de Souza. --- Goiânia, 2026.

\vspace{0.8cm}
\noindent Pré-projeto de Trabalho de Conclusão de Curso (Bacharelado em Sistemas de Informação) --- Instituto Federal de Educação, Ciência e Tecnologia de Goiás, Campus Goiânia, 2026.

\vspace{0.3cm}
\noindent Orientador: Prof. \textit{<título>}. \textit{<nome>}

\vspace{0.3cm}
\noindent 1. Evasão escolar. 2. EJA. 3. PROEJA. 4. IFG Goiânia. 5. Dashboard. 6. Análise de dados educacionais. I. Título.

% ===== SUMÁRIO =====
\newpage
\tableofcontents
\newpage

% ===== LISTA DE TABELAS E FIGURAS =====
\listoftables
\addcontentsline{toc}{section}{LISTA DE TABELAS}

\listoffigures
\addcontentsline{toc}{section}{LISTA DE FIGURAS}
\newpage

% ===== LISTA DE ABREVIATURAS =====
\section*{LISTA DE ABREVIATURAS}
\addcontentsline{toc}{section}{LISTA DE ABREVIATURAS}

\begin{tabular}{ll}
\textbf{AUC} & Area Under the Curve (Área sob a Curva) \\
\textbf{BI}  & Business Intelligence \\
\textbf{CAPES} & Coordenação de Aperfeiçoamento de Pessoal de Nível Superior \\
\textbf{EDA} & Exploratory Data Analysis (Análise Exploratória de Dados) \\
\textbf{EJA} & Educação de Jovens e Adultos \\
\textbf{IFG} & Instituto Federal de Goiás \\
\textbf{INEP} & Instituto Nacional de Estudos e Pesquisas Educacionais Anísio Teixeira \\
\textbf{KPI} & Key Performance Indicator (Indicador-Chave de Desempenho) \\
\textbf{LDB} & Lei de Diretrizes e Bases da Educação Nacional \\
\textbf{LGPD} & Lei Geral de Proteção de Dados Pessoais \\
\textbf{MEC} & Ministério da Educação \\
\textbf{PNAES} & Programa Nacional de Assistência Estudantil \\
\textbf{PNE} & Plano Nacional de Educação \\
\textbf{PROEJA} & Programa de Integração da Educação Profissional com a Educação Básica na Modalidade de Educação de Jovens e Adultos \\
\textbf{ROC} & Receiver Operating Characteristic \\
\textbf{SI} & Sistemas de Informação \\
\end{tabular}
\newpage

% ============================================================
% 1 INTRODUÇÃO
% ============================================================
\section{INTRODUÇÃO}

A Educação de Jovens e Adultos (EJA) constitui uma modalidade de ensino fundamental para a promoção da inclusão social e do resgate da cidadania de sujeitos historicamente marginalizados do processo educacional brasileiro. Conforme dados do Anuário Brasileiro da Educação Básica \cite{todos2025}, o número de matrículas na EJA caiu 34\% entre 2013 e 2023, passando de 3,6 milhões para 2,4 milhões — o menor registro da série histórica. A queda foi ainda mais acentuada no Ensino Médio, que recuou de 1,22 milhão para 888 mil matrículas no mesmo período. Paralelamente, a taxa de analfabetismo funcional atinge 17\% da população com 15 anos ou mais, enquanto 9 milhões de jovens de 15 a 29 anos (18,5\% da faixa etária) não estudam e não concluíram a Educação Básica, percentual que salta para 28\% entre os 25\% mais pobres da população \cite{todos2025}.

No âmbito da Rede Federal de Educação Profissional e Tecnológica, o Programa de Integração da Educação Profissional com a Educação Básica na Modalidade de Educação de Jovens e Adultos (PROEJA) foi instituído pelo Decreto nº 5.478/2005 \cite{mec2005}, posteriormente ampliado pelo Decreto nº 5.840/2006 \cite{mec2006}. O programa representa uma tentativa de superar a histórica separação entre educação geral e formação profissional no Brasil \cite{frignotto2005}, ofertando educação profissional técnica de nível médio integrada ao ensino médio na modalidade EJA.

O cenário nacional de declínio contrasta com a realidade de Goiás, um dos únicos três estados brasileiros que registraram aumento no total de matrículas em EJA entre 2013 e 2023, com crescimento de 1,26\% no período, saltando de 59.808 para 60.559 matrículas \cite{todos2025}. No entanto, no que tange à Educação Profissional e Tecnológica (EPT) de nível médio, Goiás apresenta apenas 6,7\% de suas matrículas nessa modalidade, contra 13,1\% da média nacional, evidenciando um potencial de expansão significativo \cite{todos2025}. No âmbito da rede federal, o estado conta com 3,4\% das matrículas de ensino médio, contexto no qual se insere o IFG Campus Goiânia.

No Instituto Federal de Goiás (IFG) -- Campus Goiânia, o PROEJA teve início com os Cursos Técnicos Integrados em Serviços de Alimentação e Cozinha, enfrentando desde sua origem desafios significativos relacionados à evasão escolar. Estudos específicos como os de \citeonline{pereira2011} e \citeonline{castro2016} revelam um percurso contraditório na construção do direito à educação, marcado por resistência institucional, focalização do programa e inadequação metodológica.

Diante deste cenário, identifica-se a necessidade de investigar de forma sistemática os padrões de evasão no PROEJA do IFG Campus Goiânia, integrando análise de dados institucionais e referencial teórico sobre educação de jovens e adultos. O desafio reside em como a tecnologia da informação, por meio de ferramentas de \textit{Business Intelligence} (BI) e análise preditiva, pode auxiliar na identificação precoce de alunos em risco de abandono e no monitoramento contínuo dos indicadores de permanência.

Nesse contexto, assume-se que o desenvolvimento de um \textit{dashboard} interativo para visualização e análise dos dados de evasão, combinado com a identificação de padrões sociodemográficos e institucionais, permitirá à gestão acadêmica tomar decisões mais informadas para a elaboração de políticas de permanência estudantil. Espera-se que esta pesquisa contribua para a redução dos índices de abandono e para o fortalecimento do PROEJA como política pública de inclusão educacional.

% ============================================================
% 2 JUSTIFICATIVA
% ============================================================
\section{JUSTIFICATIVA}

A escolha deste tema justifica-se pela relevância social e acadêmica da EJA no Brasil, modalidade que atende majoritariamente alunos pretos e pardos (76,8\% das matrículas), trabalhadores de baixa renda e pessoas em situação de vulnerabilidade social \cite{todos2025}. O perfil dos alunos da EJA se sobrepõe ao dos 9 milhões de jovens de 15 a 29 anos que não estudam e não concluíram a Educação Básica, dos quais 28\% estão entre os 25\% mais pobres da população \cite{todos2025}. A evasão escolar no PROEJA representa não apenas o desperdício de recursos públicos, mas também a perpetuação de um ciclo de exclusão educacional que compromete a mobilidade social desses indivíduos.

Ademais, o contexto goiano apresenta particularidades que reforçam a pertinência desta pesquisa. Goiás foi um dos únicos três estados brasileiros a registrar crescimento no número de matrículas em EJA entre 2013 e 2023 (+1,26\%), mas apresenta apenas 6,7\% de seus alunos de ensino médio matriculados em Educação Profissional e Tecnológica, metade da média nacional de 13,1\% \cite{todos2025}. Essa dualidade — crescimento da demanda por EJA combinado com baixa inserção em EPT — torna o IFG Campus Goiânia um lócus privilegiado para investigar os desafios e as potencialidades do PROEJA.

O presente trabalho contribui diretamente para a formação esperada no Perfil do Egresso do curso de Sistemas de Informação, uma vez que exige a aplicação integrada de conhecimentos multidisciplinares das seguintes áreas:

\begin{itemize}[leftmargin=1.25cm]
    \item[\textbf{Banco de Dados:}] modelagem e estruturação dos dados institucionais de matrícula, frequência e evasão, aplicando técnicas de modelagem dimensional para suporte à decisão;
    \item[\textbf{Business Intelligence:}] desenvolvimento de \textit{dashboard} interativo utilizando ferramentas como Power BI ou Google Looker Studio para análise de indicadores educacionais;
    \item[\textbf{Mineração de Dados:}] aplicação de técnicas de análise exploratória e identificação de padrões nos dados de evasão, permitindo a construção de modelos preditivos de risco;
    \item[\textbf{Engenharia de Software:}] levantamento de requisitos e modelagem do sistema de monitoramento, garantindo que a solução siga padrões de qualidade e usabilidade.
\end{itemize}

Do ponto de vista institucional, a pesquisa oferece ao IFG Campus Goiânia um instrumento concreto para o monitoramento da evasão, permitindo a identificação precoce de alunos em risco e a avaliação do impacto das políticas de assistência estudantil. Conforme \citeonline{oliveira2021}, a temática da evasão no PROEJA ainda carece de estudos aprofundados que integrem análise quantitativa de dados e propostas de intervenção baseadas em evidências.

A inovação deste trabalho reside na integração de técnicas de análise de dados educacionais com a proposição de medidas práticas de intervenção, consolidando o papel do Bacharel em Sistemas de Informação como um agente capaz de projetar, implementar e gerenciar soluções tecnológicas que impactam positivamente a sociedade e a gestão educacional pública.

% ============================================================
% 3 OBJETIVOS
% ============================================================
\section{OBJETIVOS}

\subsection{Objetivo Geral}

Analisar os fatores associados à evasão no PROEJA do IFG Campus Goiânia, identificando padrões sociodemográficos e institucionais, propondo medidas de intervenção e desenvolvendo um \textit{dashboard} interativo para monitoramento contínuo dos indicadores de permanência estudantil.

\subsection{Objetivos Específicos}

\begin{enumerate}[leftmargin=1.25cm]
    \item Levantar o referencial teórico sobre evasão escolar na EJA e no PROEJA, com ênfase em estudos realizados no IFG Campus Goiânia;
    \item Identificar os principais fatores que contribuem para a evasão nos cursos PROEJA do IFG Campus Goiânia, a partir da análise de dados institucionais e da literatura existente;
    \item Mapear o perfil sociodemográfico dos alunos evadidos (faixa etária, sexo, raça, renda, situação de trabalho, presença de filhos) para identificação de grupos de risco;
    \item Elaborar um conjunto de propostas de intervenção institucional, pedagógica e de apoio ao estudante visando a redução das taxas de evasão;
    \item Modelar e desenvolver um \textit{dashboard} interativo para visualização dos indicadores de evasão, utilizando ferramentas de \textit{Business Intelligence};
    \item Avaliar a eficácia do \textit{dashboard} como ferramenta de apoio à gestão acadêmica na tomada de decisões relacionadas à permanência estudantil.
\end{enumerate}

% ============================================================
% 4 REFERENCIAL TEÓRICO
% ============================================================
\section{REFERENCIAL TEÓRICO}

\subsection{Fundamentos da Educação de Jovens e Adultos no Brasil}

A Educação de Jovens e Adultos (EJA) no Brasil constitui uma modalidade educativa amparada pela Lei de Diretrizes e Bases da Educação Nacional (LDB 9.394/96), que em seu art. 37 estabelece que a educação de jovens e adultos será destinada àqueles que não tiveram acesso ou oportunidade de estudos no ensino fundamental e médio na idade própria. A LDB determina ainda que o poder público deve viabilizar e estimular o acesso e a permanência do trabalhador na escola, mediante ações integradas e complementares.

O percurso histórico da EJA no Brasil revela uma trajetória marcada pela descontinuidade de políticas públicas e pela oscilação entre diferentes concepções pedagógicas. \citeonline{dipierro2005} oferece uma análise detalhada desse processo, identificando três fases distintas: a primeira, de caráter compensatório, emerge na década de 1940 com as Campanhas de Educação de Adultos, visando erradicar o analfabetismo; a segunda, nos anos 1960, é influenciada pelo pensamento de Paulo Freire e pelos Movimentos de Cultura Popular, que propunham uma educação crítica e politizadora; a terceira, a partir da redemocratização dos anos 1980, busca consolidar a EJA como direito subjetivo, culminando na inclusão da modalidade na LDB de 1996.

A autora argumenta que a EJA oscila entre a função de ``reparação'' e a de ``educação permanente'', sendo frequentemente relegada a segundo plano nas políticas educacionais brasileiras. Essa ambivalência se reflete na alocação de recursos, na formação docente e na própria organização curricular, que muitas vezes reproduz modelos pensados para o ensino regular sem considerar as especificidades dos alunos trabalhadores.

O pensamento de \citeonline{freire2023} é fundamental para compreender a EJA sob uma perspectiva crítica e emancipadora. Em ``Pedagogia do Oprimido'', Freire propõe uma educação dialógica, que parte da realidade vivida pelos educandos, valorizando seus saberes prévios e experiências de vida. Para Freire, a educação de adultos não pode ser uma mera transferência de conhecimento, mas um ato de conhecimento onde educador e educando aprendem juntos, na problematização do mundo. Essa abordagem contrapõe-se ao que o autor chama de ``educação bancária'', na qual o aluno é tratado como um receptáculo vazio a ser preenchido pelo professor, em uma relação vertical e autoritária.

A aplicação do pensamento freireano à EJA implica reconhecer que os educandos adultos trazem consigo uma rica bagagem de conhecimentos construídos ao longo da vida, especialmente no mundo do trabalho. Ignorar esses saberes significa reproduzir a exclusão que os levou a abandonar a escola na infância ou adolescência. \citeonline{arroyo2017} aprofunda essa discussão ao caracterizar os ``sujeitos da EJA'' como ``passageiros da noite'' — trabalhadores que, após uma jornada exaustiva de trabalho, buscam na escola noturna o direito a uma educação que lhes foi negado. Arroyo argumenta que esses sujeitos são marcados por trajetórias de interrupção, trabalho precoce, responsabilidades familiares e violência institucional, sendo necessário que a escola se reorganize para acolhê-los em suas especificidades.

O Plano Nacional de Educação (PNE 2014--2024) estabeleceu metas específicas para a EJA, em particular as Metas 8, 9 e 10, que visam, respectivamente: elevar a escolaridade média da população, reduzir o analfabetismo funcional e oferecer educação profissional integrada à EJA. Contudo, dados do Anuário Brasileiro da Educação Básica \cite{todos2025} indicam que o acesso à EJA caiu 34,5\% na última década, registrando 2,4 milhões de matrículas em 2024 — o menor número da série histórica. Esse declínio evidencia o descumprimento das metas do PNE e a necessidade de políticas mais efetivas para a modalidade.

O perfil demográfico dos alunos da EJA revela a interseção entre desigualdade racial, econômica e educacional no Brasil. O Anuário 2025 \cite{todos2025} indica que 76,8\% dos alunos da EJA são pretos ou pardos, e a maioria possui renda familiar de até 1,5 salários mínimos. A pesquisa revela ainda que 22,4\% dos jovens pretos e pardos de 15 a 29 anos não estudam nem concluíram a Educação Básica, contra 15,8\% entre brancos — uma diferença de 6,6 pontos percentuais que expõe a desigualdade racial estrutural. Na zona rural, esse indicador atinge 32,2\%, o dobro da zona urbana (16\%) \cite{todos2025}.

\citeonline{silva2023} corroboram esses dados ao analisar a EJA em Goiás e Goiânia, constatando que a modalidade concentra os grupos historicamente mais vulneráveis: negros, trabalhadores informais, moradores da periferia e pessoas com defasagem idade-série superior a dois anos. Essa realidade impõe desafios específicos à gestão educacional, que precisa considerar fatores extraescolares na formulação de políticas de permanência.

\subsection{O PROEJA como Política de Integração}

O Programa Nacional de Integração da Educação Profissional com a Educação Básica na Modalidade de Educação de Jovens e Adultos (PROEJA) foi instituído pelo Decreto nº 5.478/2005 \cite{mec2005} e posteriormente ampliado pelo Decreto nº 5.840/2006 \cite{mec2006}. O programa representa uma tentativa de superar a histórica separação entre educação geral e formação profissional no Brasil, ofertando educação profissional técnica de nível médio integrada ao ensino médio na modalidade EJA.

De acordo com \citeonline{frignotto2005}, a dualidade educacional brasileira tem raízes históricas profundas: de um lado, uma educação acadêmica e propedêutica destinada às elites; de outro, uma educação profissional aligeirada e instrumental voltada para as classes trabalhadoras. O ensino médio integrado proposto pelo PROEJA visa romper com essa dicotomia ao articular trabalho, ciência, tecnologia e cultura em um currículo único, que forma o estudante simultaneamente para o exercício profissional e para a continuidade dos estudos.

O Documento Base do PROEJA \cite{mec2006} estabelece que o programa se fundamenta em três princípios: (a) a integração entre educação geral e formação profissional; (b) a formação humana integral, que supere a visão reducionista do treinamento para o mercado de trabalho; e (c) o respeito às especificidades dos sujeitos da EJA, considerando seus saberes, tempos e ritmos de aprendizagem.

Entretanto, a implementação do PROEJA tem enfrentado inúmeros desafios. \citeonline{pereira2011} aponta que o programa foi implantado de forma compulsória nas instituições federais, sem a devida preparação institucional e sem a participação efetiva dos docentes na construção dos projetos político-pedagógicos. Essa imposição gerou resistências internas e dificuldades de adaptação curricular, comprometendo a qualidade da oferta.

\citeonline{machado2011} corrobora essa análise ao observar que o PROEJA foi concebido em um contexto de disputas políticas e ideológicas, sendo marcado por contradições entre a proposta de formação integral e a realidade da precarização do trabalho docente e da falta de infraestrutura nas instituições. A autora destaca que a formação de professores para atuar na EJA integrada à educação profissional ainda é incipiente, constituindo um dos principais entraves à efetivação do programa.

No âmbito do IFG, \citeonline{castro2016} identificaram que o PROEJA foi implementado de forma focalizada, com adesão inicial de apenas uma coordenação de área (Turismo e Hospitalidade), revelando resistência institucional ao programa. As autoras argumentam que essa resistência reflete uma visão elitista da educação profissional, que historicamente privilegiou o atendimento a alunos regulares em detrimento da modalidade EJA, reproduzindo a dualidade que o programa pretendia superar.

\subsection{Evasão Escolar na EJA e no PROEJA}

A evasão escolar é um fenômeno complexo e multidimensional, que tem sido objeto de estudo de diversas áreas do conhecimento. \citeonline{tinto1975}, em seu trabalho seminal sobre evasão no ensino superior, define o fenômeno como o processo de abandono do curso ou instituição antes da conclusão, distinguindo entre evasão institucional (quando o aluno abandona uma instituição específica) e evasão sistêmica (quando abandona o sistema educacional como um todo). Tinto propõe um modelo teórico integrativo que relaciona a decisão de evadir a fatores como: integração acadêmica e social, compromisso com os objetivos educacionais, experiências institucionais prévias e características individuais.

Embora o modelo de Tinto tenha sido originalmente desenvolvido para o contexto universitário norte-americano, seus conceitos têm sido adaptados para a realidade brasileira, especialmente no que se refere à influência de fatores externos à instituição no processo de evasão. \citeonline{ferrao2003} destacam que, em contextos de desigualdade social, os determinantes da evasão extrapolam o ambiente acadêmico, estando fortemente associados a condições socioeconômicas, necessidade de trabalho e responsabilidades familiares.

No contexto da EJA e do PROEJA, a evasão assume contornos ainda mais específicos. \citeonline{oliveira2021}, em sua revisão integrativa sobre a evasão no PROEJA, analisaram 14 artigos publicados entre 2010 e 2020 e identificaram duas grandes categorias de causas: (1) fatores relacionados ao trabalho, remuneração e benefícios sociais; e (2) fatores relacionados à formação e prática docente. A primeira categoria reflete a realidade do aluno-trabalhador, que muitas vezes precisa optar entre o estudo e a sobrevivência — a incompatibilidade entre horários de trabalho e aula, o cansaço físico decorrente da dupla jornada e a necessidade de complementar a renda familiar são os motivos mais frequentemente citados. A segunda categoria revela a inadequação das práticas pedagógicas tradicionais para atender as especificidades dos estudantes jovens e adultos: currículos infantilizados, metodologias expositivas e descontextualizadas, e falta de diálogo entre o conteúdo ensinado e a realidade dos alunos.

\citeonline{nery2020} investigaram as mediações pedagógicas em turmas de EJA e constataram que predomina uma prática de ensino alicerçada na didática tradicional (``bancária''), em detrimento de uma abordagem contextualizada e dialógica. Os autores observaram que o número de professores com formação específica para atuar na EJA é insuficiente, e que a capacitação continuada para a modalidade é rara, gerando impasses e contradições no processo de ensino-aprendizagem.

\citeonline{dutra2026} analisaram os desafios contemporâneos da EJA sob a ótica da evasão e da prática pedagógica, identificando que a falta de políticas públicas consistentes e a precarização do trabalho docente agravam o quadro de abandono escolar. Os autores propõem que as estratégias de enfrentamento da evasão devem articular ações institucionais, pedagógicas e de assistência estudantil, em uma abordagem sistêmica.

\citeonline{santos2025} destacam a permanência de metodologias pedagógicas inadequadas, a descontinuidade das políticas públicas e os efeitos do envelhecimento populacional como desafios para a EJA no século XXI. Os autores argumentam que a modalidade sofre com a instabilidade política, sendo frequentemente afetada por mudanças de governo e pela falta de continuidade dos programas educacionais.

\citeonline{behrens2019} aprofundam a discussão sobre a docência na EJA, apontando que o desafio não é apenas atrair professores para a modalidade, mas formá-los para uma prática pedagógica diferenciada, que considere as especificidades dos alunos trabalhadores. Os autores defendem a adoção de metodologias ativas e contextualizadas, que partam da realidade dos educandos e promovam a articulação entre teoria e prática.

\citeonline{tolentino2025} reforça que o fortalecimento da EJA requer o comprometimento do poder público, da sociedade civil e da comunidade acadêmica na construção de políticas permanentes e eficazes, que transcendam os ciclos eleitorais e se consolidem como direito inalienável da população jovem e adulta brasileira.

\subsection{Estudos sobre o PROEJA no IFG Campus Goiânia}

A produção acadêmica sobre o PROEJA no IFG Campus Goiânia, embora ainda limitada, oferece contribuições fundamentais para a compreensão dos fatores que influenciam a evasão nesse contexto específico. Esses estudos revelam um percurso marcado por contradições entre a proposta emancipadora do programa e as dificuldades concretas de sua implementação.

\citeonline{pereira2011} desenvolveu um estudo de caso sobre o PROEJA no IFG Campus Goiânia, investigando os fatores de acesso e permanência na escola. A pesquisa, realizada com alunos dos cursos Técnicos Integrados em Serviços de Alimentação e Cozinha, revelou que os estudantes são majoritariamente trabalhadores, com idade entre 18 e 40 anos, que buscam no programa tanto a conclusão da educação básica quanto a qualificação profissional. Os principais fatores de evasão identificados foram: incompatibilidade entre horário de trabalho e aula, cansaço físico decorrente da jornada laboral, dificuldades de aprendizagem em disciplinas básicas (especialmente Matemática e Português) e falta de apoio familiar. Pereira também observou que a maioria dos alunos reside em bairros periféricos de Goiânia, enfrentando dificuldades de deslocamento e transporte público precário.

\citeonline{castro2016} analisaram o percurso contraditório do PROEJA no IFG Campus Goiânia entre 2006 e 2014, apontando que o programa foi implementado de forma focalizada, com adesão inicial de apenas uma coordenação de área (Turismo e Hospitalidade). As autoras argumentam que a resistência institucional ao PROEJA reflete uma visão elitista da educação profissional no IFG, que historicamente priorizou cursos de nível superior e técnicos regulares em detrimento da modalidade EJA. O estudo revela que a falta de planejamento institucional e a ausência de um projeto pedagógico específico para o programa comprometeram sua efetividade desde a origem.

Em sua dissertação, \citeonline{castro2018} aprofunda a análise das contradições, limites e perspectivas do PROEJA no IFG Campus Goiânia. A pesquisa examina os cursos Técnicos Integrados em Serviços de Alimentação e Cozinha e constata que, embora o programa tenha representado um avanço na democratização do acesso à educação profissional, sua implementação foi marcada por: (a) falta de materiais didáticos específicos para a EJA; (b) desvalorização dos profissionais que atuam na modalidade; (c) currículos que não dialogam adequadamente com a realidade dos alunos trabalhadores; e (d) ausência de um sistema de monitoramento e avaliação da permanência estudantil.

Conjuntamente, esses estudos indicam que o PROEJA no IFG Campus Goiânia enfrenta desafios que vão além da sala de aula, envolvendo questões estruturais, pedagógicas e institucionais. A presente pesquisa pretende dar continuidade a essa linha de investigação, incorporando a análise de dados institucionais e o desenvolvimento de ferramentas tecnológicas de monitoramento como subsídio para a formulação de políticas de permanência.

\subsection{Análise de Dados Educacionais e Mineração de Dados}

A análise de dados educacionais, também conhecida como \textit{Educational Data Mining} (EDM), tem se consolidado como uma área promissora para a compreensão dos fenômenos educacionais. \citeonline{baker2014} definem a EDM como o campo que desenvolve métodos para explorar conjuntos de dados educacionais, com o objetivo de compreender melhor os estudantes e os contextos em que aprendem. As técnicas mais comuns incluem predição, \textit{clustering}, mineração de regras de associação e análise de redes sociais.

No contexto brasileiro, \citeonline{assis2017} aplicou técnicas de mineração de dados para criar perfis de estudantes com propensão à evasão no ensino superior, utilizando algoritmos de classificação como CART (\textit{Classification and Regression Trees}) e \textit{Random Forest}. O estudo identificou que as principais características dos alunos com maior risco de evasão incluem: ingressar no primeiro semestre, possuir vínculos com mais de uma instituição e ter notas acima da média no ENEM. Assis obteve uma acurácia de aproximadamente 75\% na predição, demonstrando o potencial dessas técnicas para a identificação precoce de alunos em risco.

A modelagem preditiva de evasão tem evoluído significativamente com o avanço dos algoritmos de aprendizado de máquina. A Regressão Logística é frequentemente utilizada como modelo baseline por sua interpretabilidade, permitindo identificar o peso de cada variável na probabilidade de evasão. Já algoritmos como \textit{Random Forest} e \textit{Gradient Boosting} oferecem maior capacidade preditiva ao capturar relações não-lineares e interações entre variáveis \cite{hastie2009}. A validação cruzada com k-fold é empregada para evitar \textit{overfitting} e garantir a generalização dos modelos.

A utilização de \textit{dashboards} para monitoramento educacional tem se expandido como ferramenta de apoio à gestão acadêmica. Segundo \citeonline{few2006}, \textit{dashboards} são displays visuais das informações mais importantes necessárias para atingir um ou mais objetivos, consolidadas em uma única tela para monitoramento rápido. \citeonline{tufte2001} complementa essa visão ao defender que a visualização de dados deve priorizar a densidade de informação útil, eliminando elementos decorativos que não agregam valor à comunicação.

No contexto educacional, \textit{dashboards} de evasão permitem que gestores acompanhem em tempo real os indicadores de permanência, identifiquem precocemente alunos em risco e avaliem o impacto de políticas de assistência estudantil. \citeonline{mckinney2018} destaca que a combinação de ferramentas de análise exploratória, como as bibliotecas pandas e matplotlib em Python, com plataformas de BI como Power BI e Looker Studio, oferece um ecossistema robusto para o tratamento, análise e visualização de dados educacionais.

A presente pesquisa se insere nesse campo ao propor a integração entre análise exploratória de dados, modelagem preditiva e desenvolvimento de \textit{dashboard} para o monitoramento da evasão no PROEJA do IFG Campus Goiânia, buscando contribuir tanto para a literatura sobre EDM aplicada à EJA quanto para a prática institucional.

\begin{table}[h!]
\centering
\caption{Síntese do Referencial Teórico}
\label{tab:referencial}
\begin{tabular}{p{4cm}p{2.5cm}p{7.5cm}}
\toprule
\textbf{Autor(es)} & \textbf{Ano} & \textbf{Contribuição Principal} \\
\midrule
Freire & 2023 & Educação dialógica e emancipadora de adultos \\
Di Pierro & 2005 & Trajetória histórica e políticas públicas da EJA \\
Arroyo & 2017 & Sujeitos da EJA como ``passageiros da noite'' \\
Frigotto, Ciavatta \& Ramos & 2005 & Ensino médio integrado e superação da dualidade educacional \\
Pereira & 2011 & Fatores de acesso/permanência no PROEJA IFG \\
Castro \& Vitorette & 2016 & Contradições na implantação do PROEJA IFG \\
Oliveira \& Carmo & 2021 & Revisão integrativa: causas de evasão no PROEJA \\
Nery, Pereira \& Amaral & 2020 & Mediações pedagógicas e prática ``bancária'' na EJA \\
Machado & 2011 & Contradições e resistências na implementação do PROEJA \\
Dutra, Melo \& Silva & 2026 & Evasão e prática pedagógica na EJA contemporânea \\
Assis & 2017 & Mineração de dados para perfil de evasão \\
Baker \& Inventado & 2014 & Fundamentos da mineração de dados educacionais \\
Few & 2006 & Princípios de design de dashboards \\
Silva, Silva \& Oliveira & 2023 & EJA em Goiás e Goiânia: perfil e desafios \\
\bottomrule
\end{tabular}
\fonte{Elaborado pelo autor (2026).}
\end{table}

% ============================================================
% 5 METODOLOGIA
% ============================================================
\section{METODOLOGIA DA PESQUISA}

\subsection{Classificação da Pesquisa}

Quanto à sua natureza, esta pesquisa classifica-se como \textbf{aplicada}, pois visa gerar conhecimentos para aplicação prática na solução de um problema real -- a evasão escolar no PROEJA do IFG Campus Goiânia \cite{prodanov2013}. Do ponto de vista dos objetivos, a pesquisa é \textbf{exploratória e descritiva}: exploratória por proporcionar maior familiaridade com o problema, envolvendo levantamento bibliográfico e análise de dados; descritiva por descrever as características do fenômeno e estabelecer relações entre variáveis \cite{gil2008}.

Quanto à abordagem, caracteriza-se como \textbf{mista} \cite{creswell2010}, combinando métodos quantitativos e qualitativos:
\begin{itemize}
    \item \textbf{Quantitativa}: na análise estatística dos dados de evasão, na aplicação de testes de hipótese e na construção de modelos preditivos;
    \item \textbf{Qualitativa}: na interpretação dos fatores contextuais e institucionais que influenciam a evasão.
\end{itemize}

\subsection{Procedimentos Metodológicos}

O desenvolvimento da pesquisa será organizado em seis fases interligadas, conforme apresentado na Tabela~\ref{tab:fases}.

\begin{table}[h!]
\centering
\caption{Fases da Pesquisa}
\label{tab:fases}
\begin{tabular}{c p{3cm} p{7cm}}
\toprule
\textbf{Fase} & \textbf{Atividade} & \textbf{Descrição} \\
\midrule
1 & Revisão Sistemática & Levantamento em Google Acadêmico, SciELO, CAPES, Repositórios UnB e PUC Goiás. Descritores: ``evasão escolar'', ``PROEJA'', ``EJA'', ``IFG''. \\
2 & Coleta de Dados & Solicitação ao Registro Acadêmico do IFG: matrículas, notas, frequência, dados sociodemográficos. Anonimização conforme LGPD. \\
3 & Análise Exploratória (EDA) & Python (pandas, numpy, matplotlib, seaborn, scipy). Estatísticas descritivas, visualizações, testes de hipótese ($\chi^2$, $t$, ANOVA). \\
4 & Modelagem Preditiva & Regressão Logística e \textit{Random Forest}. Validação cruzada ($k=5$). Métricas: acurácia, precisão, revocação, F1, AUC-ROC. \\
5 & Desenvolvimento do Dashboard & Power BI ou Looker Studio: KPIs, série histórica, perfil, alerta precoce, motivos. \\
6 & Propostas de Intervenção & Recomendações em 3 eixos: institucional, pedagógico e tecnológico. \\
\bottomrule
\end{tabular}
\fonte{Elaborado pelo autor (2026).}
\end{table}

\subsection{Instrumentos de Coleta e Análise}

\begin{table}[h!]
\centering
\caption{Instrumentos da Pesquisa}
\label{tab:instrumentos}
\begin{tabular}{p{4cm} p{5.5cm} p{4.5cm}}
\toprule
\textbf{Instrumento} & \textbf{Finalidade} & \textbf{Fonte} \\
\midrule
Sistema Acadêmico (Q-Acadêmico/SUAP) & Dados de matrícula, notas, frequência & Coordenação de Registro Acadêmico \\
Questionário de perfil & Dados sociodemográficos complementares & Alunos matriculados \\
Entrevista semiestruturada & Motivos de evasão (dados qualitativos) & Alunos evadidos \\
Análise documental & Projeto Pedagógico dos cursos, relatórios & Coordenação dos cursos PROEJA \\
\bottomrule
\end{tabular}
\fonte{Elaborado pelo autor (2026).}
\end{table}

\subsection{Aspectos Éticos}

A pesquisa será submetida ao Comitê de Ética em Pesquisa (CEP) do IFG para aprovação, em conformidade com a Resolução CNS nº 510/2016. Será garantido o anonimato dos participantes e a possibilidade de desistência a qualquer momento. Os dados institucionais serão tratados conforme a Lei Geral de Proteção de Dados Pessoais (Lei nº 13.709/2018), com acesso restrito aos pesquisadores e utilização exclusiva para fins acadêmicos.

% ============================================================
% 6 CRONOGRAMA
% ============================================================
\section{CRONOGRAMA}

\begin{table}[h!]
\centering
\caption{Cronograma de Execução}
\label{tab:cronograma}
\begin{tabular}{lcccccccc}
\toprule
\textbf{Atividade} & \textbf{Abr} & \textbf{Mai} & \textbf{Jun} & \textbf{Jul} & \textbf{Ago} & \textbf{Set} & \textbf{Out} & \textbf{Nov} \\
\midrule
1. Revisão Bibliográfica     & X & X &   &   &   &   &   &   \\
2. Coleta e Preparação dos Dados &   & X & X &   &   &   &   &   \\
3. Escrita TCC 1             & X & X & X & X &   &   &   &   \\
4. EDA e Modelagem          &   &   &   & X & X &   &   &   \\
5. Desenvolvimento Dashboard &   &   &   &   & X & X &   &   \\
6. Propostas de Intervenção &   &   &   &   &   & X &   &   \\
7. Escrita TCC 2             &   &   &   &   &   & X & X & X \\
8. Preparação para Defesa   &   &   &   &   &   &   &   & X \\
\bottomrule
\end{tabular}
\fonte{Elaborado pelo autor (2026).}
\end{table}

% ============================================================
% REFERÊNCIAS
% ============================================================
\newpage
\addcontentsline{toc}{section}{REFERÊNCIAS BIBLIOGRÁFICAS}
\printbibliography[title={REFERÊNCIAS BIBLIOGRÁFICAS}]

% ============================================================
% APÊNDICE - Resultados do EDA
% ============================================================
\newpage
\begin{appendix}
\section{APÊNDICE A -- RESULTADOS DA ANÁLISE EXPLORATÓRIA DE DADOS}

Este apêndice apresenta os principais resultados obtidos na análise exploratória de dados realizada com um dataset sintético de 500 alunos, representando o perfil típico dos estudantes do PROEJA no IFG Campus Goiânia.

\subsection{Distribuição dos Alunos}

\IfFileExists{figuras/eda_evasao_perfil.png}{%
\begin{figure}[h!]
\centering
\includegraphics[width=0.85\textwidth]{figuras/eda_evasao_perfil.png}
\caption{Taxa de evasão por perfil sociodemográfico.}
\label{fig:perfil}
\end{figure}
}{}

\subsection{Distribuição dos Motivos de Evasão}

\IfFileExists{figuras/eda_motivos_evasao.png}{%
\begin{figure}[h!]
\centering
\includegraphics[width=0.65\textwidth]{figuras/eda_motivos_evasao.png}
\caption{Motivos de evasão no PROEJA IFG Goiânia.}
\label{fig:motivos}
\end{figure}
}{}

\subsection{Matriz de Correlação}

\IfFileExists{figuras/eda_matriz_correlacao.png}{%
\begin{figure}[h!]
\centering
\includegraphics[width=0.75\textwidth]{figuras/eda_matriz_correlacao.png}
\caption{Matriz de correlação dos fatores associados à evasão.}
\label{fig:correlacao}
\end{figure}
}{}

\subsection{Curva ROC do Modelo Preditivo}

\IfFileExists{figuras/eda_curva_roc.png}{%
\begin{figure}[h!]
\centering
\includegraphics[width=0.65\textwidth]{figuras/eda_curva_roc.png}
\caption{Curva ROC do modelo de regressão logística para predição de evasão.}
\label{fig:roc}
\end{figure}
}{}

\subsection{Síntese dos Resultados}

\begin{table}[h!]
\centering
\caption{Síntese dos resultados do EDA}
\label{tab:resumo_eda}
\begin{tabular}{p{5cm} p{8cm}}
\toprule
\textbf{Indicador} & \textbf{Resultado} \\
\midrule
Taxa de evasão geral & 36,4\% \\
Principal motivo & Trabalho (40,1\%) \\
Perfil de maior risco & Jovem (18-24), trabalha, com filhos, baixa renda, preto/pardo \\
Janela crítica & 1º semestre (41\% das evasões) \\
Curso com maior evasão & Cozinha (42,0\%) \\
Fatores significativos ($p<0,05$) & Trabalho, Raça, Renda, Idade \\
AUC-ROC do modelo & 0,712 \\
\bottomrule
\end{tabular}
\fonte{Elaborado pelo autor (2026).}
\end{table}

\end{appendix}

\end{document}
